\subsection{Weak Convergence in \texorpdfstring{$L^p$}{Lp}}

Weak convergence in $L^p$ keeps only the scalar information detected by conjugate functions. For $1<p<+\infty$, reflexivity turns boundedness into weak subsequential compactness; at $p=1$, bounded mass can still concentrate and destroy weak compactness unless uniform integrability is imposed.

\begin{theorem} \label{thm:lp-weak-convergence-dual-tests} \hlt{$L^p$ Weak Convergence by Dual Tests}\\
Let $(X,\mathcal M,\mu)$ be a measure space, let $1<p<+\infty$, and let $q$ be conjugate to $p$. For real $L^p(\mu)$, $(f_n)\subseteq L^p(\mu)$ and $f\in L^p(\mu)$ satisfy
$f_n\rightharpoonup f$ in $L^p(\mu)$ iff $\forall g\in L^q(\mu)$, $\int_X f_ng\,d\mu\to\int_X fg\,d\mu$.
\begin{proof}
The continuous linear functionals on real $L^p(\mu)$ are exactly the integral pairings against functions in $L^q(\mu)$ by Theorem \ref{thm:riesz-representation-lp}. Under Definition \ref{def:weak-and-weak-star-convergence}, this representation is precisely the asserted integral test.\\
\hlt{Structural Commentary}: Substitute the $L^q$ representation of every dual functional into the definition of weak convergence.\\
\hlt{Key Insight}: Riesz representation turns abstract weak tests into integral tests.
\end{proof}
\end{theorem}

\begin{proposition} \label{prop:lp-weak-limits-bounded-lsc} \hlt{Weak Limits and Norm Lower Semicontinuity}\\
Let $1<p<+\infty$. If $f_n\rightharpoonup f$ in $L^p(\mu)$, then $(f_n)$ is bounded in $L^p(\mu)$, weak sequential limits are unique, and
$\norm{f}_p\leq\liminf\limits_{n\to\infty}\norm{f_n}_p.$
\begin{proof}
Boundedness and norm lower semicontinuity follow from Proposition \ref{prop:weakly-convergent-sequences-bounded}, applied to the normed space $L^p(\mu)$. If $f_n$ converges weakly both to $f$ and $h$, then $\forall g\in L^q(\mu)$, $\int_X(f-h)g\,d\mu=0$. By the dual norm-attainment test in Proposition \ref{prop:holder-dual-norm-attainment}, $\norm{f-h}_p=0$.\\
\hlt{Structural Commentary}: General weak boundedness and lower semicontinuity are combined with an $L^q$ norming test for uniqueness.\\
\hlt{Key Insight}: Weak limits are detected by every conjugate test, including a norming test.
\end{proof}
\end{proposition}

\begin{theorem} \label{thm:lp-reflexive-sigma-finite} \hlt{Reflexivity of $L^p$, $1<p<+\infty$}\\
Let $(X,\mathcal M,\mu)$ be a $\sigma$-finite measure space and let $1<p<+\infty$. Then real $L^p(\mu)$ is reflexive.
\begin{proof}
Let $q$ be conjugate to $p$. For conjugate exponents $r,s\in(1,+\infty)$, write
$T_r(u)(v)\defeq\int_Xuv\,d\mu,\; u\in L^r(\mu),\ v\in L^s(\mu).$
Theorem \ref{thm:riesz-representation-lp} says that each $T_r:L^r(\mu)\to(L^s(\mu))^*$ is an isometric isomorphism.

Let $S\in(L^p(\mu))^{**}$. Since $T_q:L^q(\mu)\to(L^p(\mu))^*$ is bounded, $S\circ T_q$ is a bounded linear functional on $L^q(\mu)$. Hence there is $f\in L^p(\mu)$ such that
$S(T_qg)=\int_Xfg\,d\mu\;(g\in L^q(\mu)).$
If $\Phi\in(L^p(\mu))^*$, choose $g\in L^q(\mu)$ with $\Phi=T_qg$. Then
$S(\Phi)=S(T_qg)=\int_Xfg\,d\mu=\Phi(f)=Jf(\Phi).$
Thus the natural embedding $J:L^p(\mu)\to(L^p(\mu))^{**}$ is onto, and $L^p(\mu)$ is reflexive.\\
\hlt{Structural Commentary}: Apply the $L^p$--$L^q$ representation twice and identify the bidual action with evaluation.\\
\hlt{Key Insight}: The dualities $(L^p)^*=L^q$ and $(L^q)^*=L^p$ identify every bidual functional with evaluation at an $L^p$ function.
\end{proof}
\end{theorem}

\begin{theorem} \label{thm:lp-weak-sequential-compactness-bounded-sequences} \hlt{Riesz Weak Compactness Theorem}\\
Let $(X,\mathcal M,\mu)$ be a $\sigma$-finite measure space, let $1<p<+\infty$, and let $(f_n)$ be bounded in real $L^p(\mu)$. Then some subsequence converges weakly in $L^p(\mu)$ to a function $f\in L^p(\mu)$. Equivalently, for the conjugate exponent $q$,
$\int_X f_{n_k}g\,d\mu\to\int_Xfg\,d\mu \;\forall g\in L^q(\mu).$
\begin{proof}[Proof idea]
$L^p(\mu)$ is reflexive by Theorem \ref{thm:lp-reflexive-sigma-finite}. The general reflexive-space weak subsequence theorem, equivalently the sequential form supplied by Eberlein--Smulian for weakly compact bounded balls, gives a weakly convergent subsequence. The integral formulation is exactly Theorem \ref{thm:lp-weak-convergence-dual-tests}.\\
\hlt{Structural Commentary}: Reflexivity supplies weak subsequential compactness, and dual representation translates it to integral convergence.\\
\hlt{Key Insight}: Reflexivity turns bounded $L^p$ sequences into weakly compact objects.
\end{proof}
\end{theorem}

\begin{definition} \label{def:weak-sequential-compactness} \hlt{Weak Sequential Compactness}\\
A subset $K$ of a normed space $\mathcal{X}$ is \hlt{weakly sequentially compact} if every sequence in $K$ has a subsequence that converges weakly to a point of $K$.
\end{definition}

\begin{corollary} \label{cor:lp-closed-unit-ball-weakly-sequentially-compact} \hlt{The $L^p$ Unit Ball Is Weakly Sequentially Compact}\\
Let $(X,\mathcal M,\mu)$ be $\sigma$-finite and let $1<p<+\infty$. Then
$\{f\in L^p(\mu):\norm{f}_p\leq1\}$
is weakly sequentially compact in $L^p(\mu)$.
\begin{proof}
Given a sequence in the unit ball, choose a weakly convergent subsequence $f_{n_k}\rightharpoonup f$ by Theorem \ref{thm:lp-weak-sequential-compactness-bounded-sequences}. By Proposition \ref{prop:lp-weak-limits-bounded-lsc},
$\norm{f}_p\leq\liminf\limits_{k\to\infty}\norm{f_{n_k}}_p\leq1,$
so the weak limit remains in the unit ball.
\hlt{Structural Commentary}: Select a weak limit from boundedness and retain the norm constraint by lower semicontinuity.\\
\hlt{Key Insight}: Weak lower semicontinuity keeps the limit inside the closed unit ball.
\end{proof}
\end{corollary}

\begin{theorem} \label{thm:radon-riesz-lp} \hlt{Radon--Riesz Principle}\\
Let $1<p<+\infty$. In real $L^p(\mu)$, weak convergence together with convergence of norms is equivalent to strong convergence:
$f_n\rightharpoonup f$ and $\norm{f_n}_p\to\norm{f}_p$ iff $\norm{f_n-f}_p\to0$.
\hlt{Structural Commentary}: Uniform convexity of $L^p$ upgrades weak convergence plus norm convergence to strong convergence.\\
\hlt{Key Insight}: Weak convergence fixes direction; norm convergence prevents residual mass.
\end{theorem}

\begin{remark} \label{rem:radon-riesz-l2-proof-identity} \hlt{$L^2$ Proof of Radon--Riesz}\\
For $p=2$, the nontrivial direction follows from
$\norm{f_n-f}_2^2 =\int_X\abs{f_n}^2\,d\mu -2\int_Xf_nf\,d\mu +\int_X\abs f^2\,d\mu.$
Weak convergence gives $\int_Xf_nf\,d\mu\to\int_X\abs f^2\,d\mu$, and norm convergence supplies the remaining term.\\
\hlt{Key Insight}: In Hilbert space, weak convergence controls the cross term.
\end{remark}

\begin{corollary} \label{cor:radon-riesz-strong-subsequence-criterion} \hlt{Strong Subsequence Criterion}\\
Let $1<p<+\infty$ and suppose $f_n\rightharpoonup f$ in real $L^p(\mu)$. Then $(f_n)$ has a subsequence converging strongly to $f$ if and only if
$\norm{f}_p=\liminf\limits_{n\to\infty}\norm{f_n}_p.$
\begin{proof}
If $f_{n_k}\to f$ strongly, then $\norm{f_{n_k}}_p\to\norm{f}_p$, and the displayed equality follows from Proposition \ref{prop:lp-weak-limits-bounded-lsc}. Conversely, choose $n_k$ with $\norm{f_{n_k}}_p\to\liminf\limits_n\norm{f_n}_p=\norm{f}_p$. The subsequence is still weakly convergent to $f$, so $\norm{f_{n_k}-f}_p\to0$ by Theorem \ref{thm:radon-riesz-lp}.\\
\hlt{Structural Commentary}: Choose a subsequence realizing the norm liminf and invoke Radon--Riesz.\\
\hlt{Key Insight}: Once weak convergence is fixed, the only obstruction to a strongly convergent subsequence is loss of norm.
\end{proof}
\end{corollary}

\begin{example} \label{ex:l1-concentration-no-weak-subsequence} \hlt{A Bounded $L^1$ Sequence with No Weakly Convergent Subsequence}\\
On $[0,1]$, set $I_n=[0,\frac{1}{n}]$ and $f_n=n\mathbf{1}_{I_n}$. Then $\forall n,\ \norm{f_n}_1=1$. This bounded sequence has no weakly convergent subsequence in $L^1[0,1]$.
\begin{proof}
Suppose $f_{n_k}\rightharpoonup f$ in $L^1[0,1]$. If $0<c<d\leq1$, then integration against $\mathbf{1}_{[c,d]}$ is a bounded linear functional on $L^1[0,1]$, and $f_{n_k}$ vanishes on $[c,d]$ for sufficiently large $k$. Hence
$\int_c^d f\,dm=\lim\limits_{k\to\infty}\int_c^d f_{n_k}\,dm=0.$
Thus $f=0$ a.e. on $(0,1]$. Testing instead against $\mathbf{1}_{[0,1]}$ gives
$0=\int_0^1f\,dm=\lim\limits_{k\to\infty}\int_0^1f_{n_k}\,dm=1,$
a contradiction.\\
\hlt{Key Insight}: The mass stays equal to $1$ but concentrates into intervals whose lengths shrink to $0$.
\end{proof}
\end{example}

\begin{figure}[H]
  \centering
  \begin{tikzpicture}[
  x=6.9cm,
  y=0.26cm,
  >=Stealth,
  every node/.style={font=\scriptsize},
  axis/.style={black!45, ->, line width=0.45pt},
  tick/.style={black!35, line width=0.35pt},
  spikea/.style={fill=blue!15, draw=blue!70!black, line width=0.65pt},
  spikeb/.style={fill=green!12, draw=green!45!black, line width=0.65pt},
  spikec/.style={fill=orange!18, draw=orange!80!black, line width=0.65pt},
  guide/.style={black!30, dashed, line width=0.35pt}
]
  \pgfmathsetmacro{\xEight}{1/8}
  \pgfmathsetmacro{\xFour}{1/4}
  \pgfmathsetmacro{\xTwo}{1/2}

  \draw[axis] (-0.03,0) -- (1.05,0) node[right] {$x$};
  \draw[axis] (0,-0.15) -- (0,8.75) node[above] {$f_n(x)$};

  \filldraw[spikea] (0,0) rectangle (\xTwo,2);
  \filldraw[spikeb] (0,0) rectangle (\xFour,4);
  \filldraw[spikec] (0,0) rectangle (\xEight,8);

  \draw[guide] (\xTwo,0) -- (\xTwo,2.18);
  \draw[guide] (\xFour,0) -- (\xFour,4.18);
  \draw[guide] (\xEight,0) -- (\xEight,8.18);

  \foreach \x/\label in {0/$0$,\xEight/$\frac{1}{8}$,\xFour/$\frac{1}{4}$,\xTwo/$\frac{1}{2}$,1/$1$} {
    \draw[tick] (\x,0) -- (\x,-0.16);
    \node[below=2pt] at (\x,0) {\label};
  }

  \node[blue!70!black, anchor=west] at (\xTwo,2.36) {$n=2$};
  \node[green!45!black, anchor=west] at (\xFour,4.34) {$n=4$};
  \node[orange!80!black, anchor=west] at (\xEight,8.32) {$n=8$};
\end{tikzpicture}

  \caption{The functions $f_n=n\mathbf{1}_{[0,\frac{1}{n}]}$ keep total $L^1$ mass equal to $1$ while concentrating at the origin.}
  \label{fig:l1-concentration-spikes}
\end{figure}

\begin{remark} \label{rem:l1-concentration-and-uniform-integrability} \hlt{Endpoint Failure and Uniform Integrability}\\
The sequence in Example \ref{ex:l1-concentration-no-weak-subsequence} is not uniformly integrable. Thus the failure at $p=1$ is a concentration phenomenon, not merely a failure of boundedness.
\end{remark}

\begin{theorem} \label{thm:banach-saks-lp} \hlt{Banach--Saks Theorem}\\
Let $1<p<+\infty$. If $f_n\rightharpoonup f$ in real $L^p(\mu)$, then there is a subsequence $(f_{n_k})$ whose arithmetic means converge strongly to $f$ in $L^p(\mu)$:
$\left\|\frac{f_{n_1}+\cdots+f_{n_k}}{k}-f\right\|_p\to0.$
\hlt{Structural Commentary}: Select a subsequence whose successive averages exploit the uniform convexity of $L^p$.\\
\hlt{Key Insight}: Convex averaging cancels weak oscillation while preserving the weak limit.
\end{theorem}

\begin{remark} \label{rem:banach-saks-l2-selection-estimate} \hlt{The $L^2$ Selection Estimate}\\
For $p=2$, the proof may be seen directly. After replacing $f_n$ by $f_n-f$, assume $f_n\rightharpoonup0$ and $\norm{f_n}_2^2\leq M$. Choose $n_1=1$. If $S_k=f_{n_1}+\cdots+f_{n_k}$ has been chosen, weak convergence lets one choose $n_{k+1}>n_k$ so that $\abs{\int_X S_kf_{n_{k+1}}\,d\mu}\leq1$. Then
$\norm{S_{k+1}}_2^2 =\norm{S_k}^2_2+2\int_X S_kf_{n_{k+1}}\,d\mu+\norm{f_{n_{k+1}}}_2^2 \leq \norm{S_k}_2^2+2+M.$
Induction gives $\norm{\frac{S_k}{k}}_2^2\leq\frac{2+M}{k}$, so the arithmetic means converge strongly to $0$.\\
\hlt{Key Insight}: Weak convergence makes later terms nearly orthogonal to the current partial sum.
\end{remark}

\begin{definition} \label{def:closed-convex-sets-and-convex-functionals} \hlt{Closed Convex Sets and Convex Functionals}\\
Let $C$ be a subset of a normed space $\mathcal{X}$.
\begin{enumerate}[label=(\roman*)]
\setlength{\itemsep}{0pt}
\item $C$ is \hlt{convex} if $\lambda f+(1-\lambda)g\in C$ whenever $f,g\in C$ and $0\leq\lambda\leq1$.
\item $C$ is \hlt{closed} if every strongly convergent sequence in $C$ has its limit in $C$.
\item A real-valued functional $T:C\to\R$ is \hlt{continuous} if $f_n\to f$ strongly in $C$ implies $T(f_n)\to T(f)$.
\item If $C$ is convex, then $T$ is \hlt{convex} if
$T(\lambda f+(1-\lambda)g)\leq\lambda T(f)+(1-\lambda)T(g)$
$\forall f,g\in C$ and $0\leq\lambda\leq1$.
\end{enumerate}
\end{definition}

\begin{example} \label{ex:lp-closed-convex-basic-sets} \hlt{Basic Closed Convex Sets in $L^p$}\\
Let $E\subseteq\R$ be measurable and $1\leq p<+\infty$.
\begin{enumerate}[label=(\roman*)]
\setlength{\itemsep}{0pt}
\item If $g\in L^p(E)$ and $g\geq0$, then
$C_g\defeq\{f\in L^p(E):\abs{f}\leq g\text{ a.e. on }E\}$
is closed and convex.
\item The closed unit ball $\{f\in L^p(E):\norm{f}_p\leq1\}$ is closed and convex.
\end{enumerate}
\begin{proof}
Convexity of $C_g$ follows from $\abs{\lambda f+(1-\lambda)h}\leq\lambda g+(1-\lambda)g=g$. If $f_n\in C_g$ and $f_n\to f$ in $L^p(E)$, choose a subsequence converging to $f$ a.e. by Theorem \ref{thm:riesz-fischer-lp-banach}; hence $\abs f\leq g$ a.e. The unit ball is convex by Minkowski's inequality, Theorem \ref{thm:minkowski-inequality-lp}, and closed by the reverse triangle inequality.\\
\hlt{Key Insight}: An a.e.-convergent subsequence passes pointwise bounds to the limit.
\end{proof}
\end{example}

\begin{lemma} \label{lem:banach-saks-weak-limit-closed-convex-lsc} \hlt{Weak Limits and Convex Lower Semicontinuity}\\
Let $E\subseteq\R$ be measurable, let $1<p<+\infty$, let $C\subseteq L^p(E)$ be closed, bounded, and convex, and let $T:C\to\R$ be continuous and convex. If $(f_n)\subseteq C$ and $f_n\rightharpoonup f$ in $L^p(E)$, then $f\in C$ and
$T(f)\leq\liminf\limits_{n\to\infty}T(f_n).$
\begin{proof}
By Theorem \ref{thm:banach-saks-lp}, a subsequence of arithmetic means of $(f_n)$ converges strongly to $f$. These means belong to $C$ by convexity, so $f\in C$ by closedness.

Choose a subsequence, still denoted $(f_{n_k})$, such that $T(f_{n_k})\to\alpha=\liminf\limits_nT(f_n)$. Passing to a further Banach--Saks subsequence if needed, its means $a_k=\frac{f_{n_1}+\cdots+f_{n_k}}{k}$ converge strongly to $f$. Continuity and convexity give
$T(f)=\lim\limits_{k\to\infty}T(a_k) \leq\lim\limits_{k\to\infty}\frac{T(f_{n_1})+\cdots+T(f_{n_k})}{k} =\alpha.$
\hlt{Structural Commentary}: Banach--Saks places strong convex averages in $C$, and convexity transfers the liminf bound.\\
\hlt{Key Insight}: Banach--Saks converts weak convergence into strong convergence of convex averages.
\end{proof}
\end{lemma}

\begin{theorem} \label{thm:direct-method-convex-functional-lp} \hlt{Direct Method for Convex Functionals on $L^p$}\\
Let $E\subseteq\R$ be measurable and $1<p<+\infty$. If $C\subseteq L^p(E)$ is nonempty, closed, bounded, and convex, and $T:C\to\R$ is continuous and convex, then $T$ attains a minimum on $C$.
\begin{proof}
First $T$ is bounded below on $C$. Otherwise choose $f_n\in C$ with $T(f_n)\to-\infty$. By Theorem \ref{thm:lp-weak-sequential-compactness-bounded-sequences}, a subsequence converges weakly to some $f\in L^p(E)$, and $f\in C$ with $T(f)\leq\liminf T(f_n)=-\infty$ by Lemma \ref{lem:banach-saks-weak-limit-closed-convex-lsc}, impossible because $T$ is real-valued.

Let $c=\inf\limits_{h\in C}T(h)$ and choose $f_n\in C$ with $T(f_n)\to c$. Again pass to a weakly convergent subsequence with limit $f_0$. The lemma gives $f_0\in C$ and
$T(f_0)\leq\liminf\limits_{n\to\infty}T(f_n)=c.$
Thus $T(f_0)=c$.\\
\hlt{Structural Commentary}: Extract a weakly convergent minimizing subsequence and apply convex lower semicontinuity at its limit.\\
\hlt{Key Insight}: Weak sequential compactness supplies a candidate limit, and convex lower semicontinuity prevents the infimum from escaping.
\end{proof}
\end{theorem}

\begin{corollary} \label{cor:integral-convex-functional-unit-ball-minimizer} \hlt{Integral Convex Functionals Have Unit-Ball Minimizers}\\
Let $E\subseteq\R$ have finite measure, let $1<p<+\infty$, and let $\varphi:\R\to\R$ be continuous and convex. Suppose there are constants $c_1,c_2\geq0$ such that
$\abs{\varphi(s)}\leq c_1+c_2\abs{s}^p\;(s\in\R).$
Then there is $f_0\in L^p(E)$ with $\norm{f_0}_p\leq1$ such that
$\int_E\varphi(f_0)\,dm = \min\limits_{\substack{f\in L^p(E)\\ \norm{f}_p\leq1}} \int_E\varphi(f)\,dm.$
\begin{proof}
The growth bound makes $\varphi\circ f$ integrable $\forall f\in L^p(E)$, so define $T(f)=\int_E\varphi(f)\,dm$ on the closed unit ball. The ball is closed, bounded, and convex by Example \ref{ex:lp-closed-convex-basic-sets}, and $T$ is convex because $\varphi$ is convex.

It remains to note continuity. If $f_n\to f$ in $L^p(E)$, then every subsequence has a rapidly Cauchy subsequence. Along such a subsequence, the proof of Theorem \ref{thm:riesz-fischer-lp-banach} supplies an $L^p$ majorant $g$ with $\abs{f_{n_k}}\leq g$ a.e. and $f_{n_k}\to f$ a.e. Hence
$\abs{\varphi(f_{n_k})}\leq c_1+c_2g^p\in L^1(E),$
and dominated convergence gives $T(f_{n_k})\to T(f)$. This subsequence criterion proves continuity of $T$. The minimum follows from Theorem \ref{thm:direct-method-convex-functional-lp}.\\
\hlt{Structural Commentary}: A rapidly Cauchy subsequence supplies domination, after which the direct method gives a minimizer.\\
\hlt{Key Insight}: The growth bound turns strong $L^p$ convergence into dominated convergence for the integral functional.
\end{proof}
\end{corollary}
